\documentclass[10pt,a4paper]{report}
\usepackage[utf8]{inputenc}
\usepackage[spanish]{babel}
\usepackage[T1]{fontenc}
\usepackage{amsmath}
\usepackage{amsfonts}
\usepackage{amssymb}
\usepackage{amsthm}
\usepackage{stmaryrd}
\usepackage[mathscr]{euscript}
% Bibliografia
\usepackage[backend=biber]{biblatex}
\bibliography{Referencias.bib}
% Cajas
\usepackage{tcolorbox}
\usepackage{tikz}
% Comandos
\newtheorem{teo}{Teorema}
\newtheorem{prop}{Proposición}
\newtheorem{pro}{Problema}
\newcommand{\autor}{\textbf{Brayan Sandoval}}
\newcommand{\asignatura}{\textbf{Metodo de Galerkin Discontinuo}}
\newcommand{\tarea}{\textbf{Evaluación 2}}
\newcommand{\fecha}{\textbf{\today}}
\newcommand{\bs}{\boldsymbol}
\newcommand{\dx}{\textup{d}x}
\newcommand{\dt}{\textup{d}t}
\newcommand{\DG}{\textup{DG}}
\newcommand{\Th}{\mathscr{T}_{h}}
\newcommand{\Hdiv}{H\left(\textup{div};\Omega\right)}
\newcommand{\Hdivt}{H(\textup{div};\mathscr{T}_{h})}
\providecommand{\Dt}[1]{\frac{\textup{d} #1}{\textup{d}t}}
\providecommand{\Dx}[1]{\frac{\textup{d} #1}{\textup{d}x}}
\providecommand{\abs}[1]{\left\lvert#1\right\rvert}
\providecommand{\norm}[1]{\left\lVert#1\right\rVert}
\providecommand{\salto}[1]{\left\llbracket#1\right\rrbracket}
\providecommand{\prom}[1]{\left \{\!\left \{#1\right \}\!\right \}}
\providecommand{\PI}[2]{\left\langle #1,#2  \right\rangle}
\providecommand{\Pii}[2]{\left( #1,#2  \right)}
\renewcommand{\theequation}{\roman{equation}}
%texto
\usepackage{lipsum} % Genera texto aleatorio
\renewcommand*{\familydefault}{\sfdefault} % Letra mas bonita
% Figuras
\usepackage{graphicx}
% Geometría
\usepackage[left= 2 cm, right = 2 cm, top = 2 cm, bottom = 2 cm]{geometry}
\usepackage{lastpage}
% Color
\usepackage{xcolor}
\definecolor{azul}{RGB}{10,10,115}
\definecolor{amarillo}{RGB}{255,204,0}
\definecolor{rojo}{RGB}{247,0,30}
% Encabezado
\usepackage{fancyhdr}
\pagestyle{fancy}
\renewcommand{\headrulewidth}{4pt} %Aumentar grosor linea encabezado
\let\oldheadrule\headrule
\renewcommand{\headrule}{\color{azul}\oldheadrule}
\renewcommand{\footrulewidth}{4pt} %Aumentar grosor linea pie de pagina
\let\oldfootrule\footrule
\renewcommand{\footrule}{\color{azul}\oldfootrule}
\rhead{\color{azul}\autor}
\chead{\color{azul}\tarea}
\lhead{\color{azul}\asignatura}
\rfoot{\color{azul} \textbf{Pág. \thepage\ - \pageref{LastPage}}}
\cfoot{}
\lfoot{\color{azul}\fecha}
% Titulo
\title{\color{azul}\textbf{Método de Galerkin Discontinuo}\\
	\textbf{Tarea 3}}
\author{\color{azul}\autor}
\date{\color{azul}\fecha}
\begin{document}
	\section*{\underline{Introducción}}
	Considere la formulación mixta de las ecuaciones de Navier-Stokes
	\begin{align*}
		L &= \nabla \bs{u}  &&\textup{en} \ \Omega, \\
		-\nu \nabla\cdot L + \nabla\cdot\left(\bs{u}\otimes \bs{u} \right) + \nabla p &= \bs{f}  &&\textup{en}\ \Omega,\\
		\nabla\cdot \bs{u} &= 0  &&\textup{en}\ \Omega,\\
		\bs{u} &= 0   &&\textup{en}\ \partial\Omega,\\
		\int_{\Omega} p &= 0,
	\end{align*}
	donde $\bs{u}$ es la velocidad, $p$ es la presión, $\nu$ es la viscosidad cinemática y $\bs{f}\in \bs{L^{2}\left(\Omega\right)}$ es la fuerza exterior del cuerpo. El dominio $\Omega\subset\mathbb{R}^{d}$ es un poligono $(d=2)$ o un poliedro $(d=3)$. Mediante un método HDG, se tiene la siguiente formulación discreta: hallar $\left(L_{h},\bs{u}_{h},p_{h},\bs{\hat{u}}_{h}\right)\in G_{h}\times\bs{V}_{h}\times Q_{h}\times \bs{M}^{\circ}_{h}$ tal que 
	\begin{align}
		\Pii{L_{h}}{G}_{\Th} + \Pii{\bs{u}_{h}}{\nabla\cdot G}_{\Th} - \PI{\bs{\hat{u}}_{h}}{G\bs{n}}_{\partial\Th} &= 0,\label{eq:FDHDGA}\\
		\Pii{\nu L_{h}}{\nabla\bs{v}}_{\Th} - \Pii{\bs{u}_{h}\otimes\bs{\beta}}{\nabla \bs{v}}_{\Th} - \Pii{p_{h}}{\nabla\cdot\bs{v}}_{\Th} - \PI{\nu\hat{L}_{h}\bs{n}-\hat{p}_{h}\bs{n}-\left(\bs{\hat{u}}_{h}\otimes\bs{\beta}\right)\bs{n}}{\bs{v}}_{\partial\Th} &= \Pii{\bs{f}}{\bs{v}}_{\Th},\label{eq:FDHDGB}\\
		-\Pii{\bs{\hat{u}}_{h}}{\nabla q}_{\Th} + \PI{\bs{\hat{u}}_{h}\cdot \bs{n}}{q}_{\partial\Th} &= 0,\label{eq:FDHDGC}\\
		\Pii{\nu\hat{L}_{h}\bs{n} - \hat{p}_{h}\bs{n} - \left(\bs{\hat{u}}_{h}\otimes\bs{\beta}\right)\bs{n}}{\bs{\mu}}_{\partial\Th} &= 0.\label{eq:FDHDGD}
	\end{align}
	Para todo $\left(G,\bs{v},q,\bs{\mu}\right)\in G_{h}\times\bs{V}_{h}\times Q_{h}\times \bs{M}^{\circ}_{h}$. Donde,
	\begin{align*}
		\left(\nu\hat{L}_{h}-\hat{p}_{h}\right)\bs{n} &:= \nu L_{h}\bs{n} - p_{h}\bs{n} - S\left(\bs{u}_{h} - \bs{\hat{u}}_{h}\right) &&\textup{en}\ \partial\Th,\\
		S &:= S_{\bs{\beta}} + S_{\bs{n}},\\
		S_{\bs{\beta}} &:= \max\left(\bs{\beta}\cdot\bs{n},0\right)\textup{Id},\\
		S_{\bs{n}} &:= \zeta_{n} h^{-1}_{K}\bs{n}\otimes\bs{n}  &&\textup{en}\ \partial\Th,
	\end{align*} 
    donde el parámetro de estabilización $\zeta_{n}$ es un número positivo, y $\bs{\beta} = \mathbb{P}(\bs{u}_{h},\bs{\hat{u}}_{h})$ definido en \cite{CAC2017}. Ademas, los espacios donde se define la formulación discreta son:
    \begin{align*}
    	G_{h} &:= \{ G\in L^{2}\left(\Omega\right): \ G|_{K}\in P_{k}\left(K\right), \ \forall K\in \Th\},\\
    	\bs{V}_{h} &:= \{ \bs{v}\in \bs{L^{2}}\left(\Omega\right): \ \bs{v}|_{K}\in \bs{P}_{k}\left(K\right), \ \forall K\in \Th\},\\
    	Q_{h} &:= \{ p\in L^{2}_{0}\left(\Omega\right): \ p|_{K}\in P_{k}\left(K\right), \ \forall K\in \Th\},\\
    	\bs{M}_{h} &:= \{\bs{\mu}\in \bs{L^{2}}\left(\mathscr{E}_{h}\right):\ \bs{\mu}|_{F}\in \bs{P}_{k}\left(F\right), \ \forall F\in \mathscr{E}_{h}\},\\
    	\bs{M}^{\circ}_{h} &:= \{\bs{\mu}\in \bs{M}_{h}: \ \bs{\mu}|_{\partial\Omega} = 0\}.
    \end{align*}
    Con el objetivo de tener una ecuación mas compacta, se definen las formas bilineales
    \begin{align}\label{eq:FDC}
    	A_{h}\left((\bs{v},\bs{\mu}),G\right) &:= \Pii{\bs{v}}{\nabla\cdot G}_{\Th} - \PI{\bs{\mu}}{G\bs{n}}_{\partial\Th},\\
    	B_{h}\left((\bs{v},\bs{\mu}),q\right) &:= -\Pii{\bs{v}}{\nabla q}_{\Th} + \PI{q}{\bs{\mu}\cdot\bs{n}}_{\partial\Th},\\
    	J_{n}\left((\bs{v}_{1},\bs{\mu}_{1}),(\bs{v}_{2},\bs{\mu}_{2})\right) &:= \sum_{K\in\Th}\zeta_{n}h^{-1}_{K}\PI{(\bs{v}_{1}-\bs{\mu}_{1})\cdot\bs{n}}{(\bs{v}_{2}-\bs{\mu}_{2})\cdot\bs{n}}_{\partial K},\\
    	\mathcal{O}_{h}\left(\bs{\beta};(\bs{v}_{1},\bs{\mu}_{1}),(\bs{v}_{2},\bs{\mu}_{2})\right) &:= -\Pii{\bs{v}_{1}\otimes\bs{\beta}}{\nabla\bs{v}_{2}}_{\Th} + \PI{(\bs{\mu}_{1}\otimes\bs{\beta})\bs{n} + S_{\bs{\beta}}(\bs{v}_{1}-\bs{\mu}_{1})}{\bs{v}_{2}-\bs{\mu}_{2}}_{\partial\Th}.
    \end{align}
    Reemplazando esto en la formulación discreta y sumando las cuatro ecuaciones, el problema se puede reescribir como: hallar $\left(L_{h},\bs{u}_{h},p_{h},\bs{\hat{u}}_{h}\right)\in G_{h}\times\bs{V}_{h}\times Q_{h}\times \bs{M}^{\circ}_{h}$ tal que 
    \begin{align}
    	\Pii{L_{h}}{G}_{\Th} &+ A_{h}\left((\bs{u}_{h},\bs{\hat{u}}_{h}),G\right) - A_{h}\left((\bs{v},\bs{\mu}),\nu L_{h}\right) - 	B_{h}\left((\bs{v},\bs{\mu}),p_{h}\right) + B_{h}\left((\bs{u}_{h},\bs{\hat{u}}_{h}),q\right)\nonumber \\  &+ J_{n}\left((\bs{u}_{h},\bs{\hat{u}}_{h}),(\bs{v},\bs{\mu})\right) + \mathcal{O}_{h}\left(\bs{\beta};(\bs{u}_{h},\bs{\hat{u}}_{h}),(\bs{v},\bs{\mu})\right) = \Pii{\bs{f}}{\bs{v}}_{\Th}, \label{eq:FDHDGC}
    \end{align} 
   donde $\bs{\beta} = \mathbb{P}(\bs{u}_{h},\bs{\hat{u}}_{h})$, para todo $\left(G,\bs{v},q,\bs{\mu}\right)\in G_{h}\times\bs{V}_{h}\times Q_{h}\times \bs{M}^{\circ}_{h}$.
   \newpage
   \section*{\underline{Existencia y Unicidad}}
    Estamos interesados en probar la existencia y unicidad de la formulación discreta definida en la introducción, esto se podrá garantizar con el siguiente teorema
    \begin{teo}\label{teo2}
    	Si la cantidad $\nu^{-2}\norm{\bs{f}}_{\Omega}$ es lo suficientemente pequeña, el método HDG (\ref*{eq:FDHDGA}-\ref*{eq:FDHDGD}) tiene única solución $\left(L_{h},\bs{u}_{h},p_{h},\bs{\hat{u}}_{h}\right)\in G_{h}\times\bs{V}_{h}\times Q_{h}\times \bs{M}^{\circ}_{h}$, y ademas
    	\begin{align*}
    		||| \left(\bs{u}_{h},\bs{\hat{u}}_{h}\right) |||_{1,h} \leq C\nu^{-1}\norm{\bs{f}}_{\Omega}
    	\end{align*}
        para una constante $C$ independiente de $\nu$, la discretización y la solución exacta.
    \end{teo}
    \begin{proof}
    El objetivo de este trabajo será mostrar este teorema usando el teorema del punto fijo de Brower, para esto debemos definir un operador que este definido en espacios los cuales sean compactos y que ademas sea una contracción. Se propone el operador $\mathcal{F}:\bs{Z}_{h}\longrightarrow \bs{Z}_{h}$, donde 
    \begin{align*}
    	\bs{Z}_{h} := \{ (\bs{v},\bs{\mu})\in \bs{V}_{h}\times \bs{M}^{\circ}_{h}:\ B_{h}\left((\bs{v},\bs{\mu}),q\right) = 0,\ \forall q\in Q_{h} \}.
    \end{align*}
    El operador $\mathcal{F}$ toma elementos $(\bs{w},\bs{\hat{w}})\in Z_{h}$ y lo transforma en un $(\bs{u},\bs{\hat{u}})\in Z_{h}$ el cual es parte de la solución $(L_{\bs{w}},\bs{u},p_{\bs{w}},\bs{\hat{u}})\in  G_{h}\times\bs{V}_{h}\times Q_{h}\times \bs{M}^{\circ}_{h}$ de 
    \begin{align*}
    	\Pii{L_{\bs{w}}}{G}_{\Th} &+ A_{h}\left((\bs{u},\bs{\hat{u}}),G\right) - A_{h}\left((\bs{v},\bs{\mu}),\nu L_{\bs{w}}\right) - 	B_{h}\left((\bs{v},\bs{\mu}),p_{\bs{w}}\right) + B_{h}\left((\bs{u},\bs{\hat{u}}),q\right)\\ &+ J_{n}\left((\bs{u},\bs{\hat{u}}),(\bs{v},\bs{\mu})\right) + \mathcal{O}_{h}\left(\mathbb{P}(\bs{w},\bs{\hat{w}});(\bs{u},\bs{\hat{u}}),(\bs{v},\bs{\mu})\right) = \Pii{\bs{f}}{\bs{v}}_{\Th},
    \end{align*} 
    para todo $\left(G,\bs{v},q,\bs{\mu}\right)\in G_{h}\times\bs{V}_{h}\times Q_{h}\times \bs{M}^{\circ}_{h}$. Notemos que si se tiene $\mathcal{F}(\bs{u}_{h},\bs{\hat{u}}_{h}) = (\bs{u}_{h},\bs{\hat{u}}_{h})$ se obtiene la existencia de la solución discreta del método HDG, si este punto fijo es único se tiene la unicidad. Antes de probar las hipótesis del teorema del punto fijo veamos que $\bs{\mathcal{F}}$ \textbf{esta bien definido}, esto es, el operador es biyectivo, para esto usaremos la alternativa de Fredholm, esto lo podemos hacer ya que el operador se reduce a un sistema de ecuaciones cuadrado. Dado $(\bs{w},\bs{\hat{w}}) := (\bs{0},\bs{0})$ y $\bs{f} := \bs{0}$, se define $\mathcal{F}(\bs{0},\bs{0}) = (\bs{u}_{\bs{0}},\bs{\hat{u}}_{\bs{0}})$ tal que $(L_{\bs{0}},\bs{u}_{\bs{0}},p_{\bs{0}},\bs{\hat{u}}_{\bs{0}})\in  G_{h}\times\bs{V}_{h}\times Q_{h}\times \bs{M}^{\circ}_{h}$ soluciona el sistema
    \begin{align}
    	\Pii{L_{\bs{0}}}{G}_{\Th} &+ A_{h}\left((\bs{u}_{\bs{0}},\bs{\hat{u}}_{\bs{0}}),G\right) - A_{h}\left((\bs{v},\bs{\mu}),\nu L_{\bs{0}}\right) - 	B_{h}\left((\bs{v},\bs{\mu}),p_{\bs{0}}\right) + B_{h}\left((\bs{u}_{\bs{0}},\bs{\hat{u}}_{\bs{0}}),q\right)\nonumber\\ &+ J_{n}\left((\bs{u}_{\bs{0}},\bs{\hat{u}}_{\bs{0}}),(\bs{v},\bs{\mu})\right) + \mathcal{O}_{h}\left(\mathbb{P}(\bs{0},\bs{0});(\bs{u}_{\bs{0}},\bs{\hat{u}}_{\bs{0}}),(\bs{v},\bs{\mu})\right) = 0,\label{eq:fcero}
    \end{align} 
    para todo $\left(G,\bs{v},q,\bs{\mu}\right)\in G_{h}\times\bs{V}_{h}\times Q_{h}\times \bs{M}^{\circ}_{h}$. En particular, para $G = \nu L_{\bs{0}},\ \bs{v} = \bs{u}_{\bs{0}},\ q = p_{\bs{0}}$ y $\bs{\mu} = \bs{\hat{u}}_{\bs{0}}$ se obtiene 
    \begin{align*}
    	\abs{\nu}\norm{L_{\bs{0}}}^{2}_{\Th} + J_{n}\left((\bs{u}_{\bs{0}},\bs{\hat{u}}_{\bs{0}}),(\bs{u}_{\bs{0}},\bs{\hat{u}}_{\bs{0}})\right) = 0,
    \end{align*}
    usando la definición de $J_{n}$ se obtiene
    \begin{align*}
    	\abs{\nu}\norm{L_{\bs{0}}}^{2}_{\Th} + \sum_{K\in\Th} \zeta_{n}h^{-1}_{K}\norm{\left(\bs{u}_{\bs{0}} - \bs{\hat{u}}_{\bs{0}}\right)\cdot\bs{n}}^{2}_{\partial K} = 0
    \end{align*}
    De aquí se tiene que $L_{\bs{0}} = 0$ en $\Omega$ y $\bs{u}_{\bs{0}} = \bs{\hat{u}}_{\bs{0}}$ en $\partial \Th$. Como $\bs{\hat{u}}_{\bs{0}}\in \bs{M}^{\circ}_{h}$, se tiene que $\bs{u}_{\bs{0}} = \bs{\hat{u}}_{\bs{0}} = \bs{0}$ en $\partial\Omega$, ademas, dado $e\in \mathscr{E}^{i}_{h}$
    \begin{align*}
    	\salto{\bs{u}_{\bs{0}}} = \bs{u}_{\bs{0}}^{+}\cdot \bs{n}^{+} +  \bs{u}_{\bs{0}}^{-}\cdot \bs{n}^{-} = \bs{\hat{u}}_{\bs{0}}\cdot \bs{n}^{+} + \bs{\hat{u}}_{\bs{0}}\cdot \bs{n}^{-} = \bs{\hat{u}}_{\bs{0}}\cdot(\bs{n}^{+} + \bs{n}^{-}) = 0,
    \end{align*}
    es decir, $\bs{u}_{\bs{0}}$ es continua en $\Omega$. Reemplazando lo obtenido en \ref{eq:fcero} se obtiene
    \begin{align*}
    	A_{h}\left((\bs{u}_{\bs{0}},\bs{\hat{u}}_{\bs{0}}),G\right) - 	B_{h}\left((\bs{v},\bs{\mu}),p_{\bs{0}}\right) + B_{h}\left((\bs{u}_{\bs{0}},\bs{\hat{u}}_{\bs{0}}),q\right) + J_{n}\left((\bs{u}_{\bs{0}},\bs{\hat{u}}_{\bs{0}}),(\bs{v},\bs{\mu})\right)  = 0,
    \end{align*}
    esto se puede reescribir como sigue
    \begin{align*}
    	\Pii{\bs{u}_{\bs{0}}}{\nabla\cdot G}_{\Th} &- \PI{\bs{\hat{u}}_{\bs{0}}}{G\bs{n}}_{\partial \Th} + \Pii{\bs{v}}{\nabla p_{\bs{0}}}_{\Th} - \PI{p_{\bs{0}}}{\bs{\mu}\cdot\bs{n}}_{\partial\Th} - \Pii{\bs{u}_{\bs{0}}}{\nabla q}_{\Th} + \PI{q}{\bs{\hat{u}}_{\bs{0}}\cdot\bs{n}}_{\partial\Th}\\ &+ \sum_{K\in\Th} \zeta_{n}h^{-1}_{K}\PI{\left(\bs{u}_{\bs{0}} - \bs{\hat{u}}_{\bs{0}}\right)\cdot\bs{n}}{\left(\bs{v}-\bs{\mu}\right)\cdot\bs{n}}_{\partial K} = 0,
    \end{align*}
    dado que $\bs{u}_{\bs{0}} = \bs{\hat{u}}_{\bs{0}}$ en $\partial \Th$ se obtiene
    \begin{align*}
    	\Pii{\bs{u}_{\bs{0}}}{\nabla\cdot G}_{\Th} - \PI{\bs{\hat{u}}_{\bs{0}}}{G\bs{n}}_{\partial \Th} + \Pii{\bs{v}}{\nabla p_{\bs{0}}}_{\Th} - \PI{p_{\bs{0}}}{\bs{\mu}\cdot\bs{n}}_{\partial\Th} - \Pii{\bs{u}_{\bs{0}}}{\nabla q}_{\Th} + \PI{q}{\bs{\hat{u}}_{\bs{0}}\cdot\bs{n}}_{\partial\Th} = 0,
    \end{align*}
    esto se puede reescribir como 
    \begin{align*}
    	\sum_{K\in\Th}\left[  \Pii{\bs{u}_{\bs{0}}}{\nabla\cdot G}_{K} - \PI{\bs{\hat{u}}_{\bs{0}}}{G\bs{n}}_{\partial K} + \Pii{\bs{v}}{\nabla p_{\bs{0}}}_{K} - \PI{p_{\bs{0}}}{\bs{\mu}\cdot\bs{n}}_{\partial K} - \Pii{\bs{u}_{\bs{0}}}{\nabla q}_{K} + \PI{q}{\bs{\hat{u}}_{\bs{0}}\cdot\bs{n}}_{\partial K} \right] = 0.
    \end{align*}
    \newpage
    Aplicando integración por partes a $\Pii{\bs{u}_{\bs{0}}}{\nabla\cdot G}_{K}$ con $K\in \Th$ se obtiene
    \begin{align*}
    	\sum_{K\in\Th}\left[  -\Pii{G}{\nabla\bs{u}_{\bs{0}}}_{K} + \PI{G\bs{n}}{\bs{u}_{\bs{0}}}_{\partial K} - \PI{\bs{\hat{u}}_{\bs{0}}}{G\bs{n}}_{\partial K} + \Pii{\bs{v}}{\nabla p_{\bs{0}}}_{K} - \PI{p_{\bs{0}}}{\bs{\mu}\cdot\bs{n}}_{\partial K} - \Pii{\bs{u}_{\bs{0}}}{\nabla q}_{K} + \PI{q}{\bs{\hat{u}}_{\bs{0}}\cdot\bs{n}}_{\partial K} \right] = 0,
    \end{align*}
    esto se puede reescribir como
    \begin{align*}
    	\sum_{K\in\Th}\left[  -\Pii{G}{\nabla\bs{u}_{\bs{0}}}_{K} +  \PI{\bs{u}_{\bs{0}} - \bs{\hat{u}}_{\bs{0}}}{G\bs{n}}_{\partial K} + \Pii{\bs{v}}{\nabla p_{\bs{0}}}_{K} - \PI{p_{\bs{0}}}{\bs{\mu}\cdot\bs{n}}_{\partial K} - \Pii{\bs{u}_{\bs{0}}}{\nabla q}_{K} + \PI{q}{\bs{\hat{u}}_{\bs{0}}\cdot\bs{n}}_{\partial K} \right] = 0,
    \end{align*}
    dado que $\bs{u}_{\bs{0}} = \bs{\hat{u}}_{\bs{0}}$ en $\partial \Th$ se obtiene
     \begin{align*}
    	\sum_{K\in\Th}\left[  -\Pii{G}{\nabla\bs{u}_{\bs{0}}}_{K} + \Pii{\bs{v}}{\nabla p_{\bs{0}}}_{K} - \PI{p_{\bs{0}}}{\bs{\mu}\cdot\bs{n}}_{\partial K} - \Pii{\bs{u}_{\bs{0}}}{\nabla q}_{K} + \PI{q}{\bs{\hat{u}}_{\bs{0}}\cdot\bs{n}}_{\partial K} \right] = 0,
    \end{align*}
    aplicando nuevamente integración por partes, esta vez a $\Pii{\bs{v}}{\nabla p_{\bs{0}}}_{K}$ y $\Pii{\bs{u}_{\bs{0}}}{\nabla q}_{K}$ para todo $K\in \Th$ se llega a
    \begin{align*}
    		\sum_{K\in\Th}\left[  -\Pii{G}{\nabla\bs{u}_{\bs{0}}}_{K} - \Pii{\nabla\cdot \bs{v}}{p_{\bs{0}}}_{K} + \Pii{\nabla\cdot \bs{u}_{\bs{0}}}{q}_{K} + \PI{\left(\bs{\mu} - \bs{v}\right)\cdot\bs{n}}{p_{\bs{0}}}_{\partial K} \right] = 0
    \end{align*}
    para todo $\left(G,\bs{v},q,\bs{\mu}\right)\in G_{h}\times\bs{V}_{h}\times Q_{h}\times \bs{M}^{\circ}_{h}$. En particular, para $G = \nabla \bs{u}_{\bs{0}}$, $\bs{v} = \bs{u}_{\bs{0}}$, $q = p_{\bs{0}}$ y $\bs{\mu} = \bs{\hat{u}}_{\bs{0}}$ se obtiene
    \begin{align*}
    	-\norm{\nabla\bs{u}_{\bs{0}}}^{2}_{\Th} - \Pii{\nabla\cdot \bs{u}_{\bs{0}}}{p_{\bs{0}}}_{\Th} + \Pii{\nabla\cdot \bs{u}_{\bs{0}}}{p_{\bs{0}}}_{\Th} + \PI{\left(\bs{\hat{u}}_{\bs{0}} - \bs{u}_{\bs{0}}\right)\cdot\bs{n}}{p_{\bs{0}}}_{\partial \Th} = 0,
    \end{align*}
    esto es 
    \begin{align*}
    	\norm{\nabla\bs{u}_{\bs{0}}}_{\Th} = 0 \iff \nabla\bs{u}_{\bs{0}} = \bs{0} \iff \bs{u}_{\bs{0}}|_{K} = \textup{Cte}_{K}, \forall K\in \Th.
    \end{align*}
    Como $\bs{u}_{\bs{0}}$ es continua y ademas se anula en $\partial\Omega$, entonces se tiene que $\bs{u}_{\bs{0}} = \bs{0}$. Nos falta probar que $p_{\bs{0}} = 0$, para esto usemos que
    \begin{align*}
    	\sum_{K\in\Th}\left[  -\Pii{G}{\nabla\bs{u}_{\bs{0}}}_{K} - \Pii{\nabla\cdot \bs{v}}{p_{\bs{0}}}_{K} + \Pii{\nabla\cdot \bs{u}_{\bs{0}}}{q}_{K} + \PI{\left(\bs{\mu} - \bs{v}\right)\cdot\bs{n}}{p_{\bs{0}}}_{\partial K} \right] = 0
    \end{align*}
    para todo $\left(G,\bs{v},q,\bs{\mu}\right)\in G_{h}\times\bs{V}_{h}\times Q_{h}\times \bs{M}^{\circ}_{h}$. Como $\nabla \bs{u}_{\bs{0}} = \bs{0}$ se tiene que 
    \begin{align*}
    	\sum_{K\in\Th}\left[   - \Pii{\nabla\cdot \bs{v}}{p_{\bs{0}}}_{K} + \PI{\left(\bs{\mu} - \bs{v}\right)\cdot\bs{n}}{p_{\bs{0}}}_{\partial K} \right] = 0
    \end{align*}
    para todo $\left(G,\bs{v},q,\bs{\mu}\right)\in G_{h}\times\bs{V}_{h}\times Q_{h}\times \bs{M}^{\circ}_{h}$. Integrando por partes se obtiene 
    \begin{align}\label{eq:pql}
    	\sum_{K\in\Th}\left[   - \Pii{\bs{v}}{\nabla p_{\bs{0}}}_{K} + \PI{\bs{v}\cdot\bs{n}}{p_{\bs{0}}}_{\partial K} + \PI{\left(\bs{\mu} - \bs{v}\right)\cdot\bs{n}}{p_{\bs{0}}}_{\partial K} \right] = 0
    \end{align}
    para todo $\left(G,\bs{v},q,\bs{\mu}\right)\in G_{h}\times\bs{V}_{h}\times Q_{h}\times \bs{M}^{\circ}_{h}$. En particular, para $\bs{\mu} = \bs{0}$ y $\bs{v} = \nabla p_{\bs{0}}$ se obtiene que
    \begin{align*}
    	\norm{\nabla p_{\bs{0}}}_{\Th} = 0 \iff  \nabla p_{\bs{0}} = \bs{0} \iff p|_{K} = \textup{Cte}_{K}, \forall K\in \Th,
    \end{align*}
    veamos que $p_{\bs{0}}$ es continuo en $\Omega$, para esto usemos \eqref{eq:pql} y el hecho que $ \nabla p_{\bs{0}} = \bs{0}$
    \begin{align*}
    	\sum_{K\in\Th}\left[ \PI{\bs{v}\cdot\bs{n}}{p_{\bs{0}}}_{\partial K} + \PI{\left(\bs{\mu} - \bs{v}\right)\cdot\bs{n}}{p_{\bs{0}}}_{\partial K} \right] = 0
    \end{align*}
    para todo $\left(G,\bs{v},q,\bs{\mu}\right)\in G_{h}\times\bs{V}_{h}\times Q_{h}\times \bs{M}^{\circ}_{h}$. Desarrollando se obtiene
    \begin{align}\label{eq:aa}
    	\sum_{K\in\Th} \PI{\bs{\mu} \cdot\bs{n}}{p_{\bs{0}}}_{\partial K}  = 0, \ \forall \mu\in \bs{M}^{\circ}_{h}
    \end{align} 
    usando la siguiente propiedad
    \begin{align*}
    	\sum_{K\in\Th} \PI{\bs{\mu} \cdot\bs{n}}{p_{\bs{0}}}_{\partial K} = \PI{\salto{\bs{\mu}}}{\prom{p_{\bs{0}}}}_{\mathscr{E}^{i}_{h}} + \PI{\prom{\bs{\mu}}}{\salto{p_{\bs{0}}}}_{\mathscr{E}^{\i}_{h}}  + \PI{\bs{\mu}\cdot\bs{n}}{p_{\bs{0}}}_{\mathscr{E}^{\partial}_{h}},  \ \forall \bs{\mu}\in \bs{M}^{\circ}_{h}.
    \end{align*}
    Usando el hecho que las funciones en $\bs{M}^{\circ}_{h}$ son continuas respecto a las aristas de la triangulación, se anulan en la frontera de $\Omega$ y usando \eqref{eq:aa} se obtiene que 
    \begin{align*}
    	\PI{\bs{\mu}}{\salto{p_{\bs{0}}}}_{\mathscr{E}^{i}_{h}} = 0,\ \forall \bs{\mu}\in \bs{M}^{\circ}_{h},
    \end{align*} 
    en particular, para $\hat{e}\in \mathscr{E}^{i}_{h}$ se puede definir
    \begin{align*}
    	\bs{\mu} = \begin{cases}
    		\salto{p_{\bs{0}}}& \text{ si } e = \hat{e} \\ 
    		0 & \text{ si } e \neq \hat{e} 
    	\end{cases}
    \end{align*}
    luego,
    \begin{align*}
    	\norm{\salto{p_{\bs{0}}}}_{\hat{e}} = 0,\ \forall \hat{e}\in \mathscr{E}^{i}_{h} \Longrightarrow \norm{\salto{p_{\bs{0}}}}_{\mathscr{E}^{i}_{h}} = 0 \iff \salto{p_{\bs{0}}} = 0\  \text{en}\ \mathscr{E}^{i}_{h}
    \end{align*}
    de aquí se deduce que $p_{\bs{0}}$ es continua en $\Omega$ y por tanto $p_{\bs{0}}(x) = \textup{Cte}$ para todo $x\in \Omega$, dado que $p_{\bs{0}}\in L^{2}_{0}(\Omega)$ y que $0<|\Omega|<+\infty$ se tiene que
    \begin{align*}
    	\int_{\Omega}p_{\bs{0}} = 0 \iff p_{\bs{0}}|\Omega| = 0 \Longrightarrow p_{\bs{0}} = 0.
    \end{align*}
    Finalmente se deduce que $(L_{\bs{0}},\bs{u}_{\bs{0}},p_{\bs{0}},\bs{\hat{u}}_{\bs{0}}) = (0,\bs{0},0,\bs{0})$, es decir, $\mathcal{F}$ es inyectivo y por alternativa de Fredholm $\mathcal{F}$ es biyectivo y por tanto esta bien definido. Ahora, mostraremos que el operador $\bs{\mathcal{F}}$ \textbf{va de} $\bs{K}_{h}$ \textbf{a} $\bs{K}_{h}$ donde 
    \begin{align*}
    	\bs{K}_{h} := \{(\bs{v},\bs{\mu})\in \bs{Z}_{h}:\ |||(\bs{v},\bs{\mu})|||_{1,h}\leq C_{2}C^{2}_{\textup{HDG}} \nu^{-1}\norm{\bs{f}}_{\Omega} \}
    \end{align*}
    por la forma del conjunto es fácil ver que es compacto y convexo. Dado $(\bs{w}_{h},\hat{\bs{w}}_{h})\in \bs{Z}_{h}$, luego existe $(\bs{u}_{h},\hat{\bs{u}}_{h})$ tal que $(L_{h},\bs{u}_{h},p_{h},\bs{\hat{u}}_{h})\in G_{h}\times\bs{V}_{h}\times Q_{h}\times \bs{M}^{\circ}_{h}$ soluciona \ref{eq:FDHDGC} con $\bs{\beta} = \mathbb{P}(\bs{w}_{h},\hat{\bs{w}}_{h})$, por \textbf{proposición 2.2} de \cite{CAC2017} se tiene 
    \begin{align*}
    	|||(\bs{u}_{h},\bs{\hat{u}}_{h})|||_{1,h} \leq C_{\textup{HDG}} \left(\norm{L_{h}}^{2}_{\Th} + \sum_{K\in\Th} h^{-1}_{K}\norm{(\bs{u}_{h} - \bs{\hat{u}}_{h})\cdot \bs{n}}^{2}_{\partial K} \right)^{\frac{1}{2}},
    \end{align*} 
    esto se puede reescribir
    \begin{align*}
    	\nu|||(\bs{u}_{h},\bs{\hat{u}}_{h})|||^{2}_{1,h} \leq \nu C^{2}_{\textup{HDG}} \left(\norm{L_{h}}^{2}_{\Th} + \sum_{K\in\Th} h^{-1}_{K}\norm{(\bs{u}_{h} - \bs{\hat{u}}_{h})\cdot \bs{n}}^{2}_{\partial K} \right),
    \end{align*} 
    ahora, reemplazando $(L_{h},\bs{u}_{h},p_{h},\bs{\hat{u}}_{h})$ en \ref{eq:FDHDGC} y considerando $G = \nu L_{h}$, $\bs{v} = \bs{u}_{h}$, $\bs{\mu} = \bs{\hat{u}}_{h}$ y $q = p_{h}$, se obtiene 
    \begin{align*}
    	\Pii{L_{h}}{\nu L_{h}}_{\Th} &+ A_{h}\left((\bs{u}_{h},\bs{\hat{u}}_{h}),\nu L_{h}\right) - A_{h}\left((\bs{u}_{h},\bs{\hat{u}}_{h}),\nu L_{h}\right) - B_{h}\left((\bs{u}_{h},\bs{\hat{u}}_{h}),p_{h}\right) + B_{h}\left((\bs{u}_{h},\bs{\hat{u}}_{h}),p_{h}\right)\\ &+ J_{n}\left((\bs{u}_{h},\bs{\hat{u}}_{h}),(\bs{u}_{h},\bs{\hat{u}}_{h})\right) + \mathcal{O}_{h}\left(\bs{\beta}; (\bs{u}_{h},\bs{\hat{u}}_{h}),(\bs{u}_{h},\bs{\hat{u}}_{h})\right) = \Pii{\bs{f}}{\bs{u}_{h}}_{\Th}, 
    \end{align*}
    esto se puede reescribir
    \begin{align*}
    	\nu\norm{L_{h}}^{2}_{\Th} + J_{n}\left((\bs{u}_{h},\bs{\hat{u}}_{h}),(\bs{u}_{h},\bs{\hat{u}}_{h})\right) + \mathcal{O}_{h}\left(\bs{\beta}; (\bs{u}_{h},\bs{\hat{u}}_{h}),(\bs{u}_{h},\bs{\hat{u}}_{h})\right) = \Pii{\bs{f}}{\bs{u}_{h}}_{\Th}
    \end{align*}
    usando la definición de $J_{n}$ se obtiene
    \begin{align*}
    	\nu\norm{L_{h}}^{2}_{\Th} + \sum_{K\in\Th}\zeta_{n}h^{-1}_{K}\norm{(\bs{u}_{h}-\bs{\hat{u}}_{h})\cdot\bs{n}}^{2}_{\partial K} + \mathcal{O}_{h}\left(\bs{\beta}; (\bs{u}_{h},\bs{\hat{u}}_{h}),(\bs{u}_{h},\bs{\hat{u}}_{h})\right) = \Pii{\bs{f}}{\bs{u}_{h}}_{\Th},
    \end{align*}
    notemos que $\bs{\beta} = \mathbb{P}(\bs{w}_{h},\hat{\bs{w}}_{h})$, donde $(\bs{w}_{h},\hat{\bs{w}}_{h})\in \bs{Z}_{h}$ por \textbf{proposición 2.1} de \cite{CAC2017} se tiene que 
    \begin{align*}
    	\mathbb{P}(\bs{w}_{h},\hat{\bs{w}}_{h})\in \Hdiv \ \textup{y}\ \nabla\cdot\mathbb{P}(\bs{w}_{h},\hat{\bs{w}}_{h}) = 0\ \textup{en}\ \Omega.
    \end{align*}
    Así, por \textbf{proposición 3.6} de \cite{CAC2017} como $(\bs{u}_{h},\bs{\hat{u}}_{h})\in \bs{V}_{h}\times\bs{M}_{h}$, entonces
    \begin{align*}
    	\mathcal{O}_{h}\left(\bs{\beta}; (\bs{u}_{h},\bs{\hat{u}}_{h}),(\bs{u}_{h},\bs{\hat{u}}_{h})\right) \geq 0
    \end{align*}
    recordemos que $\zeta_{n}$ es un termino de estabilización positivo y $\zeta_{n}\gg 1$, de esto se desprende que 
    \begin{align*}
    	\nu\norm{L_{h}}^{2}_{\Th} + \nu \sum_{K\in\Th} h^{-1}_{K}\norm{(\bs{u}_{h}-\bs{\hat{u}}_{h})\cdot\bs{n}}^{2}_{\partial K}\leq \Pii{\bs{f}}{\bs{u}_{h}}_{\Th},
    \end{align*}
    por Cauchy-Schwarz
    \begin{align*}
    	\abs{\Pii{\bs{f}}{\bs{u}_{h}}_{\Th}} \leq \norm{\bs{f}}_{\Omega}\norm{\bs{u}_{h}}_{\Th}
    \end{align*}
    juntando todo
     \begin{align*}
    	\nu ||| (\bs{u}_{h},\bs{\hat{u}}_{h}) |||^{2}_{1,h}\leq \nu C^{2}_{\textup{HDG}}\left(
    	\norm{L_{h}}^{2}_{\Th} + \sum_{K\in\Th}h^{-1}_{K}\norm{(\bs{u}_{h}-\bs{\hat{u}}_{h})\cdot\bs{n}}^{2}_{\partial K}  \right)\leq  C^{2}_{\textup{HDG}}\Pii{\bs{f}}{\bs{u}_{h}}_{\Th}\leq  C^{2}_{\textup{HDG}}\norm{\bs{f}}_{\Omega}\norm{\bs{u}_{h}}_{\Th}.
    \end{align*}
    Por \textbf{proposición A2} de \cite{CAC2017} para $q = 2$, se tiene que 
    \begin{align*}
    	\norm{\bs{u}_{h}}_{\Th} \leq C_{2}|||(\bs{u}_{h},\bs{\hat{u}}_{h}) |||_{1,h},
    \end{align*}
    ya que $(\bs{u}_{h},\bs{\hat{u}}_{h})\in V(h)\times\bs{M}^{\circ}_{h}$. Considerando esto se obtiene
    \begin{align*}
    	\nu ||| (\bs{u}_{h},\bs{\hat{u}}_{h}) |||^{2}_{1,h}\leq C^{2}_{\textup{HDG}}C_{2}\norm{\bs{f}}_{\Omega}|||(\bs{u}_{h},\bs{\hat{u}}_{h}) |||_{1,h} \iff ||| (\bs{u}_{h},\bs{\hat{u}}_{h}) |||_{1,h}\leq C^{2}_{\textup{HDG}}C_{2}\nu^{-1}\norm{\bs{f}}_{\Omega}
    \end{align*}
    luego, $\mathcal{F}(\bs{w}_{h},\hat{\bs{w}}_{h}) = (\bs{u}_{h},\bs{\hat{u}}_{h})\in \bs{K}_{h}$, por tanto usando la sobreyectividad de $\mathcal{F}$ se tiene que es un operador que va de $\bs{K}_{h}$ a $\bs{K}_{h}$. Ahora probemos que $\bs{\mathcal{F}}$ \textbf{es una contracción en} $\bs{K}_{h}$, sean $(\bs{w}_{1},\bs{\hat{w}}_{1}),(\bs{w}_{2},\bs{\hat{w}}_{2})\in \bs{K}_{h}$, luego existen $(\bs{u}_{1},\bs{\hat{u}}_{1}),(\bs{u}_{2},\bs{\hat{u}}_{2})\in \bs{K}_{h}$ tales que $(L_{1},\bs{u}_{1},p_{1},\bs{\hat{u}}_{1}),(L_{2},\bs{u}_{2},p_{2},\bs{\hat{u}}_{2})\in G_{h}\times\bs{V}_{h}\times Q_{h}\times \bs{M}^{\circ}_{h}$ solucionan \ref{eq:FDHDGC} con $\bs{\beta}_{1} = \mathbb{P}(\bs{w}_{1},\hat{\bs{w}}_{1})$ y $\bs{\beta}_{2} = \mathbb{P}(\bs{w}_{2},\hat{\bs{w}}_{2})$, si restamos ambas ecuaciones y definimos $\delta_{L} := L_{1} - L_{2}$, $\delta_{\bs{u}} := \bs{u}_{1} - \bs{u}_{2}$, $\delta_{\bs{\hat{u}}} := \bs{\hat{u}}_{1} - \bs{\hat{u}}_{2}$ y $\delta_{p} := p_{1} - p_{2}$ se obtiene
    \begin{align*}
    	\Pii{\delta_{L}}{G}_{\Th} &+ A_{h}\left((\delta_{\bs{u}},\delta_{\bs{\hat{u}}}),G\right) - A_{h}\left((\bs{v},\bs{\mu}),\nu\delta_{L}\right) - B_{h}\left((\bs{v},\bs{\mu}),\delta_{p}\right) + B_{h}\left((\delta_{\bs{u}},\delta_{\bs{\hat{u}}}),q\right)\\ &+ J_{n}\left((\delta_{\bs{u}},\delta_{\bs{\hat{u}}}),(\bs{v}, \bs{\mu})\right) + \mathcal{O}_{h}\left(\mathbb{P}(\bs{w}_{1},\bs{\hat{w}}_{1});(\bs{u}_{1},\bs{\hat{\bs{u}}}_{1}),(\bs{v},\bs{\mu})\right) - \mathcal{O}_{h}\left(\mathbb{P}(\bs{w}_{2},\bs{\hat{w}}_{2});(\bs{u}_{2},\bs{\hat{\bs{u}}}_{2}),(\bs{v},\bs{\mu})\right) = 0,
    \end{align*}
    para todo $\left(G,\bs{v},q,\bs{\mu}\right)\in G_{h}\times\bs{V}_{h}\times Q_{h}\times \bs{M}^{\circ}_{h}$. En particular para $\left(G,\bs{v},q,\bs{\mu}\right) := \left(\nu\delta_{L},\delta_{\bs{u}},\delta_{p},\delta_{\bs{\hat{u}}}\right)$ se obtiene
    \begin{align}
    	\nu \norm{\delta_{L}}^{2}_{\Th} + J_{n}\left((\delta_{\bs{u}},\delta_{\bs{\hat{u}}}),(\delta_{\bs{u}},\delta_{\bs{\hat{u}}})\right) + \mathcal{O}_{h}\left(\mathbb{P}(\bs{w}_{1},\bs{\hat{w}}_{1});(\bs{u}_{1},\bs{\hat{\bs{u}}}_{1}),(\delta_{\bs{u}},\delta_{\bs{\hat{u}}})\right) - \mathcal{O}_{h}\left(\mathbb{P}(\bs{w}_{2},\bs{\hat{w}}_{2});(\bs{u}_{2},\bs{\hat{\bs{u}}}_{2}),(\delta_{\bs{u}},\delta_{\bs{\hat{u}}})\right) = 0,\label{eq:xdd}
    \end{align}
    aplicando \textbf{proposición 2.2} de \cite{CAC2017}, se tiene
    \begin{align*}
    	\nu||| (\delta_{\bs{u}},\delta_{\bs{\hat{u}}}) |||^{2}_{1,h} \leq C^{2}_{\textup{HDG}} \nu \left(\norm{\delta_{L}}^{2}_{\Th} + \sum_{K\in\Th}h^{-1}_{K}\norm{(\delta_{\bs{u}}-\delta_{\bs{\hat{u}}})\cdot \bs{n}}^{2}_{\partial K} \right) = C^{2}_{\textup{HDG}} \left(\nu\norm{\delta_{L}}^{2}_{\Th} + \sum_{K\in\Th}\nu h^{-1}_{K}\norm{(\delta_{\bs{u}}-\delta_{\bs{\hat{u}}})\cdot \bs{n}}^{2}_{\partial K} \right) 
    \end{align*}
    usando que $\nu\leq\zeta_{n}$ se obtiene que 
    \begin{align*}
    	\nu||| (\delta_{\bs{u}},\delta_{\bs{\hat{u}}}) |||^{2}_{1,h}  \leq C^{2}_{\textup{HDG}} \left(\nu\norm{\delta_{L}}^{2}_{\Th} + J_{n}\left((\delta_{\bs{u}},\delta_{\bs{\hat{u}}}),(\delta_{\bs{u}},\delta_{\bs{\hat{u}}})\right) \right) 
    \end{align*}
    de \ref{eq:xdd} se desprende que 
    \begin{align*}
    	\nu\norm{\delta_{L}}^{2}_{\Th} + J_{n}\left((\delta_{\bs{u}},\delta_{\bs{\hat{u}}}),(\delta_{\bs{u}},\delta_{\bs{\hat{u}}})\right) = \mathcal{O}_{h}\left(\mathbb{P}(\bs{w}_{2},\bs{\hat{w}}_{2});(\bs{u}_{2},\bs{\hat{\bs{u}}}_{2}),(\delta_{\bs{u}},\delta_{\bs{\hat{u}}})\right)-\mathcal{O}_{h}\left(\mathbb{P}(\bs{w}_{1},\bs{\hat{w}}_{1});(\bs{u}_{1},\bs{\hat{\bs{u}}}_{1}),(\delta_{\bs{u}},\delta_{\bs{\hat{u}}})\right),
    \end{align*}
    luego,
    \begin{align*}
    	\nu||| (\delta_{\bs{u}},\delta_{\bs{\hat{u}}}) |||^{2}_{1,h}  &\leq C^{2}_{\textup{HDG}}\left(\mathcal{O}_{h}\left(\mathbb{P}(\bs{w}_{2},\bs{\hat{w}}_{2});(\bs{u}_{2},\bs{\hat{\bs{u}}}_{2}),(\delta_{\bs{u}},\delta_{\bs{\hat{u}}})\right)-\mathcal{O}_{h}\left(\mathbb{P}(\bs{w}_{1},\bs{\hat{w}}_{1});(\bs{u}_{1},\bs{\hat{\bs{u}}}_{1}),(\delta_{\bs{u}},\delta_{\bs{\hat{u}}})\right) \right)\\
    	&= C^{2}_{\textup{HDG}}\left(\mathcal{O}_{h}\left(\mathbb{P}(\bs{w}_{2},\bs{\hat{w}}_{2});(\bs{u}_{2},\bs{\hat{\bs{u}}}_{2}),(\delta_{\bs{u}},\delta_{\bs{\hat{u}}})\right)-\mathcal{O}_{h}\left(\mathbb{P}(\bs{w}_{1},\bs{\hat{w}}_{1});(\bs{u}_{1},\bs{\hat{\bs{u}}}_{1}),(\delta_{\bs{u}},\delta_{\bs{\hat{u}}})\right)\right.\\ &+ \left. \mathcal{O}_{h}\left(\mathbb{P}(\bs{w}_{2},\bs{\hat{w}}_{2});(\bs{u}_{1},\bs{\hat{\bs{u}}}_{1}),(\delta_{\bs{u}},\delta_{\bs{\hat{u}}})\right) - \mathcal{O}_{h}\left(\mathbb{P}(\bs{w}_{2},\bs{\hat{w}}_{2});(\bs{u}_{1},\bs{\hat{\bs{u}}}_{1}),(\delta_{\bs{u}},\delta_{\bs{\hat{u}}})\right) \right)\\
    	&= C^{2}_{\textup{HDG}}\left(\mathcal{O}_{h}\left(\mathbb{P}(\bs{w}_{2},\bs{\hat{w}}_{2});(\bs{u}_{1},\bs{\hat{\bs{u}}}_{1}),(\delta_{\bs{u}},\delta_{\bs{\hat{u}}})\right)-\mathcal{O}_{h}\left(\mathbb{P}(\bs{w}_{1},\bs{\hat{w}}_{1});(\bs{u}_{1},\bs{\hat{\bs{u}}}_{1}),(\delta_{\bs{u}},\delta_{\bs{\hat{u}}})\right)\right.\\ &+ \left. \mathcal{O}_{h}\left(\mathbb{P}(\bs{w}_{2},\bs{\hat{w}}_{2});(\bs{u}_{2},\bs{\hat{\bs{u}}}_{2}),(\delta_{\bs{u}},\delta_{\bs{\hat{u}}})\right)  - \mathcal{O}_{h}\left(\mathbb{P}(\bs{w}_{2},\bs{\hat{w}}_{2});(\bs{u}_{1},\bs{\hat{\bs{u}}}_{1}),(\delta_{\bs{u}},\delta_{\bs{\hat{u}}})\right) \right)\\
    	&= C^{2}_{\textup{HDG}}\left(\mathcal{O}_{h}\left(\mathbb{P}(\bs{w}_{2},\bs{\hat{w}}_{2});(\bs{u}_{1},\bs{\hat{\bs{u}}}_{1}),(\delta_{\bs{u}},\delta_{\bs{\hat{u}}})\right)-\mathcal{O}_{h}\left(\mathbb{P}(\bs{w}_{1},\bs{\hat{w}}_{1});(\bs{u}_{1},\bs{\hat{\bs{u}}}_{1}),(\delta_{\bs{u}},\delta_{\bs{\hat{u}}})\right)\right.\\
    	&-\left.\mathcal{O}_{h}\left(\mathbb{P}(\bs{w}_{2},\bs{\hat{w}}_{2});(\delta_{\bs{u}},\delta_{\bs{\hat{u}}}),(\delta_{\bs{u}},\delta_{\bs{\hat{u}}})\right) \right),
    \end{align*}
    Por \textbf{proposición 3.6} de \cite{CAC2017} se tiene que 
    \begin{align*}
    	\mathcal{O}_{h}\left(\mathbb{P}(\bs{w}_{2},\bs{\hat{w}}_{2});(\delta_{\bs{u}},\delta_{\bs{\hat{u}}}),(\delta_{\bs{u}},\delta_{\bs{\hat{u}}})\right)\geq 0,
    \end{align*}
    así, se deduce que
    \begin{align*}
    	\nu||| (\delta_{\bs{u}},\delta_{\bs{\hat{u}}}) |||^{2}_{1,h}  \leq C^{2}_{\textup{HDG}}\left(\mathcal{O}_{h}\left(\mathbb{P}(\bs{w}_{2},\bs{\hat{w}}_{2});(\bs{u}_{1},\bs{\hat{\bs{u}}}_{1}),(\delta_{\bs{u}},\delta_{\bs{\hat{u}}})\right)-\mathcal{O}_{h}\left(\mathbb{P}(\bs{w}_{1},\bs{\hat{w}}_{1});(\bs{u}_{1},\bs{\hat{\bs{u}}}_{1}),(\delta_{\bs{u}},\delta_{\bs{\hat{u}}})\right)\right).
    \end{align*} 
    Usando la \textbf{proposición 3.4} de \cite{CAC2017} se tiene que
    \begin{align*}
    	\nu||| (\delta_{\bs{u}},\delta_{\bs{\hat{u}}}) |||^{2}_{1,h}  &\leq C^{2}_{\textup{HDG}}C_{\mathcal{O}}\norm{\mathbb{P}(\bs{w}_{2},\bs{\hat{w}}_{2}) - \mathbb{P}(\bs{w}_{1},\bs{\hat{w}}_{1})}_{1,h}\norm{(\bs{u}_{1},\bs{\hat{\bs{u}}}_{1})}_{1,h}\norm{(\delta_{\bs{u}},\delta_{\bs{\hat{u}}})}_{1,h}\\
    	&= C^{2}_{\textup{HDG}}C_{\mathcal{O}}\norm{\mathbb{P}(\bs{w}_{2}-\bs{w}_{1},\bs{\hat{w}}_{2}-\bs{\hat{w}}_{1})}_{1,h}\norm{(\bs{u}_{1},\bs{\hat{\bs{u}}}_{1})}_{1,h}\norm{(\delta_{\bs{u}},\delta_{\bs{\hat{u}}})}_{1,h}.
    \end{align*}
    Por \textbf{proposición 2.1} de \cite{CAC2017}
    \begin{align*}
    	\nu||| (\delta_{\bs{u}},\delta_{\bs{\hat{u}}}) |||^{2}_{1,h}  \leq C^{2}_{\textup{HDG}}C_{\mathcal{O}}C_{\textup{stab},1}\norm{(\bs{w}_{2}-\bs{w}_{1},\bs{\hat{w}}_{2}-\bs{\hat{w}}_{1})}_{1,h}\norm{(\bs{u}_{1},\bs{\hat{\bs{u}}}_{1})}_{1,h}\norm{(\delta_{\bs{u}},\delta_{\bs{\hat{u}}})}_{1,h}.
    \end{align*}
    Dado que $(\bs{u}_{1},\bs{\hat{u}}_{1})\in\bs{K}_{h}$ se obtiene
    \begin{align*}
    	||| (\bs{u}_{1},\bs{\hat{u}}_{1}) |||_{1,h}\leq C^{2}_{\textup{HDG}}C_{2}\nu^{-1}\norm{\bs{f}}_{\Omega}.
    \end{align*} 
    Juntando todo se obtiene 
    \begin{align*}
    	\nu||| (\delta_{\bs{u}},\delta_{\bs{\hat{u}}}) |||^{2}_{1,h} \leq C^{4}_{\textup{HDG}}C_{2}C_{\mathcal{O}}C_{\textup{stab},1}\nu^{-1}\norm{(\bs{w}_{2}-\bs{w}_{1},\bs{\hat{w}}_{2}-\bs{\hat{w}}_{1})}_{1,h}\norm{\bs{f}}_{\Omega}\norm{(\delta_{\bs{u}},\delta_{\bs{\hat{u}}})}_{1,h},
    \end{align*}
    esto se puede reescribir como
    \begin{align*}
    	||| (\delta_{\bs{u}},\delta_{\bs{\hat{u}}}) |||_{1,h} \leq C^{4}_{\textup{HDG}}C_{2}C_{\mathcal{O}}C_{\textup{stab},1}\nu^{-2}\norm{\bs{f}}_{\Omega}\norm{(\bs{w}_{2}-\bs{w}_{1},\bs{\hat{w}}_{2}-\bs{\hat{w}}_{1})}_{1,h}.
    \end{align*}
    Por tanto, $\mathcal{F}$ es una contracción en $\bs{K}_{h}$ si se tiene que 
    \begin{align*}
    	\nu^{-2}\norm{\bs{f}}_{\Omega} < \frac{1}{C_{2}C^{4}_{\textup{HDG}}C_{\mathcal{O}}C_{\textup{stab},1}}.
    \end{align*}
    Esto es, si $\nu^{-2}\norm{\bs{f}}_{\Omega}$ es una cantidad pequeña, por teorema del punto fijo, existe un único $(\bs{u}_{h},\bs{\hat{u}}_{h})\in \bs{K}_{h}$ tal que $\mathcal{F}\left(\bs{u}_{h},\bs{\hat{u}}_{h}\right) = (\bs{u}_{h},\bs{\hat{u}}_{h})$, lo cual demuestra el teorema.
    \end{proof}
    \begin{prop}
    	Sea $\bs{\beta}\in\{\bs{v}\in\Hdiv:\ \nabla\cdot\bs{v}=0,\ \bs{v}|_{K}\in\bs{H}^{1}(K),\ \forall K\in\Th \}$. Luego,
    	\begin{align*}
    		\mathcal{O}_{h}\left(\bs{\beta};(\bs{v},\bs{\mu}),(\bs{v},\bs{\mu})\right) = \PI{(S_{\bs{\beta}}-\frac{1}{2}\bs{\beta}\cdot\bs{n}\textup{Id})(\bs{v}-\bs{\mu})}{\bs{v}-\bs{\mu}}_{\partial\Th} \geq 0 \hspace{3mm} \forall(\bs{v},\bs{\mu})\in\bs{V}_{h}\times\bs{M}_{h}.
    	\end{align*}
    \end{prop}
    \begin{proof}
    	Dado $\bs{\beta}\in\{\bs{v}\in\Hdiv:\ \nabla\cdot\bs{v}=0,\ \bs{v}|_{K}\in\bs{H}^{1}(K),\ \forall K\in\Th \}$, se tiene por integración por partes
    	\begin{align*}
    		\Pii{\nabla\cdot(\bs{v}\otimes\bs{\beta})}{\bs{v}}_{\Th} = -\Pii{(\bs{v}\otimes\bs{\beta})}{\nabla\bs{v}}_{\Th} + \PI{(\bs{v}\otimes\bs{\beta})\bs{n}}{\bs{v}}_{\partial\Th}
    	\end{align*}
        Considerando que
        \begin{align*}
        	\nabla\cdot(\bs{v}\otimes\bs{\beta}) = \begin{pmatrix}
        		\sum_{j=1}^{d}\partial_{j}(v_{1}\beta_{j})\\ 
        		\vdots\\ 
        		\sum_{j=1}^{d}\partial_{j}(v_{d}\beta_{j})
        	\end{pmatrix}
        \end{align*}
        fijando $i\in\{1,\dots,d\}$ se tiene
        \begin{align*}
        	\sum_{j=1}^{d}\partial_{j}(v_{1}\beta_{j}) = \sum_{j=1}^{d}\left[\partial_{j}(v_{i})\beta_{j} + v_{i}\partial_{j}(\beta_{j})\right] = \sum_{j=1}^{d}\partial_{j}(v_{i})\beta_{j} + v_{i}\nabla\cdot\bs{\beta} =  \sum_{j=1}^{d}\partial_{j}(v_{i})\beta_{j}
        \end{align*}
        Así
        \begin{align*}
        	\Pii{\nabla\cdot(\bs{v}\otimes\bs{\beta})}{\bs{v}}_{\Th} = \Pii{\begin{pmatrix}
        		\sum_{j=1}^{d}\partial_{j}(v_{1}\beta_{j})\\ 
        		\vdots\\ 
        		\sum_{j=1}^{d}\partial_{j}(v_{d}\beta_{j})
        	\end{pmatrix}}{\begin{pmatrix}
        	v_{1}\\ 
        	\vdots\\ 
        	v_{d}
        \end{pmatrix}}_{\Th} =  \Pii{\begin{pmatrix}
        \sum_{j=1}^{d}\partial_{j}(v_{1})\beta_{j}\\ 
        \vdots\\ 
        \sum_{j=1}^{d}\partial_{j}(v_{d})\beta_{j}
    \end{pmatrix}}{\begin{pmatrix}
    v_{1}\\ 
    \vdots\\ 
    v_{d}
\end{pmatrix}}_{\Th} = \int_{\Th}\sum_{i,j=1}^{d}\beta_{j}v_{i}\partial_{j}(v_{i})
        \end{align*}
        de forma similar
        \begin{align*}
        	\Pii{(\bs{v}\otimes\bs{\beta})}{\nabla\bs{v}}_{\Th} = \sum_{i,j=1}^{d} v_{i}\beta_{j}\partial_{j}(v_{i}) = \Pii{\nabla\cdot(\bs{v}\otimes\bs{\beta})}{\bs{v}}_{\Th},
        \end{align*}
        usando la formula de integración por partes
        \begin{align*}
        	\Pii{\nabla\cdot(\bs{v}\otimes\bs{\beta})}{\bs{v}}_{\Th} = \frac{1}{2}\PI{(\bs{v}\otimes\bs{\beta})\bs{n}}{\bs{v}}_{\Th}
        \end{align*}
        luego, por definición de $\mathcal{O}_{h}$
        \begin{align*}
        	\mathcal{O}_{h}\left(\bs{\beta};(\bs{v},\bs{\mu}),(\bs{v},\bs{\mu})\right) &= -\Pii{\bs{v}\otimes\bs{\beta}}{\nabla\bs{v}}_{\Th} + \PI{(\bs{\mu}\otimes\bs{\beta})\bs{n}+S_{\bs{\beta}}(\bs{v}-\bs{\mu})}{\bs{v}-\bs{\mu}}_{\partial\Th}\\
        	&= \Pii{\nabla\cdot(\bs{v}\otimes\bs{\beta})}{\bs{v}}_{\Th} - \PI{(\bs{v}\otimes\bs{\beta})\bs{n}}{\bs{v}}_{\partial\Th} + \PI{(\bs{\mu}\otimes\bs{\beta})\bs{n}+S_{\bs{\beta}}(\bs{v}-\bs{\mu})}{\bs{v}-\bs{\mu}}_{\partial\Th}\\
        	&= \frac{1}{2}\PI{(\bs{v}\otimes\bs{\beta})\bs{n}}{\bs{v}}_{\Th} - \PI{(\bs{v}\otimes\bs{\beta})\bs{n}}{\bs{v}}_{\partial\Th} + \PI{(\bs{\mu}\otimes\bs{\beta})\bs{n}+S_{\bs{\beta}}(\bs{v}-\bs{\mu})}{\bs{v}-\bs{\mu}}_{\partial\Th}
        \end{align*}
        Como $\bs{\mu}$ es una función univaluada y $\bs{\beta}\in\Hdiv$, se tiene que $\PI{(\bs{\mu}\otimes\bs{\beta})\bs{n}}{\bs{\mu}}_{\partial\Th} = 0$. Usando esto
        \begin{align*}
        	\mathcal{O}_{h}\left(\bs{\beta};(\bs{v},\bs{\mu}),(\bs{v},\bs{\mu})\right) &= \frac{1}{2}\PI{(\bs{v}\otimes\bs{\beta})\bs{n}}{\bs{v}}_{\Th} - \PI{(\bs{v}\otimes\bs{\beta})\bs{n}}{\bs{v}}_{\partial\Th} +\frac{1}{2}\PI{(\bs{v}\otimes\bs{\beta})\bs{n}}{\bs{\mu}}_{\partial\Th}-\frac{1}{2}\PI{(\bs{v}\otimes\bs{\beta})\bs{n}}{\bs{\mu}}_{\partial\Th}\\ &+ \PI{(\bs{\mu}\otimes\bs{\beta})\bs{n}+S_{\bs{\beta}}(\bs{v}-\bs{\mu})}{\bs{v}-\bs{\mu}}_{\partial\Th}\\
        	&= \frac{1}{2}\PI{(\bs{v}\otimes\bs{\beta})\bs{n}}{\bs{v}}_{\Th} - \PI{(\bs{v}\otimes\bs{\beta})\bs{n}}{\bs{v}}_{\partial\Th} +\frac{1}{2}\PI{(\bs{v}\otimes\bs{\beta})\bs{n}}{\bs{\mu}}_{\partial\Th}-\frac{1}{2}\PI{(\bs{v}\otimes\bs{\beta})\bs{n}}{\bs{\mu}}_{\partial\Th}\\ &+ \PI{(\bs{\mu}\otimes\bs{\beta})\bs{n}}{\bs{v}-\bs{\mu}}_{\partial\Th} +  \PI{S_{\bs{\beta}}(\bs{v}-\bs{\mu})}{\bs{v}-\bs{\mu}}_{\partial\Th}\\
        	&= \frac{1}{2}\PI{(\bs{\mu}\otimes\bs{\beta})\bs{n}}{\bs{\mu}}_{\partial\Th}-\frac{1}{2}\PI{(\bs{v}\otimes\bs{\beta})\bs{n}}{\bs{v}}_{\Th} +\frac{1}{2}\PI{(\bs{v}\otimes\bs{\beta})\bs{n}}{\bs{\mu}}_{\partial\Th} -\frac{1}{2}\PI{(\bs{v}\otimes\bs{\beta})\bs{n}}{\bs{\mu}}_{\partial\Th}\\ &+ \PI{(\bs{\mu}\otimes\bs{\beta})\bs{n}}{\bs{v}-\bs{\mu}}_{\partial\Th} +  \PI{S_{\bs{\beta}}(\bs{v}-\bs{\mu})}{\bs{v}-\bs{\mu}}_{\partial\Th}\\
        	&= -\frac{1}{2}\PI{((\bs{v}-\bs{\mu})\otimes\bs{\beta})\bs{n}}{\bs{\mu}}-\frac{1}{2}\PI{(\bs{v}\otimes\bs{\beta})\bs{n}}{\bs{v}-\bs{\mu}}_{\partial\Th} 
        	+ \PI{(\bs{\mu}\otimes\bs{\beta})\bs{n}}{\bs{v}-\bs{\mu}}_{\partial\Th} +  \PI{S_{\bs{\beta}}(\bs{v}-\bs{\mu})}{\bs{v}-\bs{\mu}}_{\partial\Th}\\
        	&= \frac{1}{2}\PI{(\bs{\mu}\otimes\bs{\beta})\bs{n}}{\bs{\mu}}_{\partial\Th}-\frac{1}{2}\PI{(\bs{v}\otimes\bs{\beta})\bs{n}}{\bs{v}}_{\Th} +\frac{1}{2}\PI{(\bs{v}\otimes\bs{\beta})\bs{n}}{\bs{\mu}}_{\partial\Th} -\frac{1}{2}\PI{(\bs{v}\otimes\bs{\beta})\bs{n}}{\bs{\mu}}_{\partial\Th}\\ &+ \PI{(\bs{\mu}\otimes\bs{\beta})\bs{n}}{\bs{v}-\bs{\mu}}_{\partial\Th} +  \PI{S_{\bs{\beta}}(\bs{v}-\bs{\mu})}{\bs{v}-\bs{\mu}}_{\partial\Th}\\
        	&= -\frac{1}{2}\PI{((\bs{v}-\bs{\mu})\otimes\bs{\beta})\bs{n}}{\bs{\mu}}-\frac{1}{2}\PI{(\bs{v}\otimes\bs{\beta})\bs{n}}{\bs{v}-\bs{\mu}}_{\partial\Th} \\
        	&+ \frac{1}{2}\PI{(\bs{\mu}\otimes\bs{\beta})\bs{n}}{\bs{v}-\bs{\mu}}_{\partial\Th} +\frac{1}{2}\PI{(\bs{\mu}\otimes\bs{\beta})\bs{n}}{\bs{v}-\bs{\mu}}_{\partial\Th} +  \PI{S_{\bs{\beta}}(\bs{v}-\bs{\mu})}{\bs{v}-\bs{\mu}}_{\partial\Th}\\
        	&= -\frac{1}{2}\PI{((\bs{v}-\bs{\mu})\otimes\bs{\beta})\bs{n}}{\bs{v}-\bs{\mu}}_{\partial\Th}+  \PI{S_{\bs{\beta}}(\bs{v}-\bs{\mu})}{\bs{v}-\bs{\mu}}_{\partial\Th}
        \end{align*}
        De aquí se sigue que
        \begin{align*}
        	\mathcal{O}_{h}\left(\bs{\beta};(\bs{v},\bs{\mu}),(\bs{v},\bs{\mu})\right) = \PI{(S_{\bs{\beta}}-\frac{1}{2}\bs{\beta}\cdot\bs{n}\textup{Id})(\bs{v}-\bs{\mu})}{\bs{v}-\bs{\mu}}_{\partial\Th}
        \end{align*}
        Y por definición de $S_{\bs{\beta}}$ se concluye que $	\mathcal{O}_{h}\left(\bs{\beta};(\bs{v},\bs{\mu}),(\bs{v},\bs{\mu})\right)\geq 0$.
    \end{proof}
    \begin{thebibliography}{0}
    	\bibitem[CAC2017]{CAC2017} Aycil Cesmelioglu, Bernardo Cockburn and Weifeng Qiu, Analysis of a hybridizable discontinuous Galerkin
    	method for the steady-state incompressible Navier-Stokes equations, Math. Comp. 86 (2017), 1643-1670.
    \end{thebibliography}
\end{document}